\documentclass[output=paper,colorlinks,citecolor=brown]{langscibook}
\ChapterDOI{10.5281/zenodo.21237011}
\author{Andrew Merritt\affiliation{Georgetown University}}
%\ORCIDs{}

\title{De-adjectival and de-``adjectival'': Greek κρύπτω `cover' and Proto-Germanic *\textit{χreuðaną} `deck, equip, adorn'}
\shorttitlerunninghead{De-adjectival and de-``adjectival''}


\abstract{This article clarifies the relationships of three secondary roots reflected in, e.g., 1) Greek κρύπτω ‘cover’ and Tocharian B \textit{kraup}- ‘gather, amass’ (: *\textit{kreu̯bʰ}-), 2) Old Church Slavonic \textit{kryti} ‘cover’ and Lithuanian \textit{kráuti}, 1sg. prs. \textit{kráuju} ‘pile up’ \mbox{(: *\textit{kreu̯h\textsubscript{x}}-),} and 3) Old English \textit{hréodan}* ‘deck, adorn’ (: *\textit{kreu̯dʰ}-), in light of our current understanding of the derivational behavior of nominal compounds formed on the basis of the predicative instrumental periphrasis. After identifying *\textit{kreh\textsubscript{2}}- ‘heap, strew, cover’ (Latvian \textit{krãju}, \textit{krât} ‘put aside, accumulate, amass’) as the primary root, the author proposes that its \textit{u}-stem derivative *\textit{kró/éh\textsubscript{2}-u}- ‘heap, covering’ was the base of an *-\textit{i̯e/o}-present *\textit{kreh\textsubscript{2}-u-i̯e/o}- ‘provide with covering’, whose reanalysis as a primary verb gave rise to *\textit{kreu̯h\textsubscript{2}}-. Since the original denominal stem  could only appear in the present system, there was need elsewhere for the use of its base’s instrumental with the light verbs *\textit{dʰeh\textsubscript{1}}- ‘put, make’ and *\textit{bʰuh\textsubscript{x}}- ‘become’. On the basis of the anticausative periphrasis *\textit{kréh\textsubscript{2}-u-h\textsubscript{1} bʰuh\textsubscript{x}-} ‘become covered’ (cf.\ Ved.\ \textit{gúhā} \textit{bhū}- ‘become hidden’) was formed an adjective *\textit{kr̥h\textsubscript{2}-u-bʰ(u̯)-ó}- ‘heaped, covered’ {>}{>} *\textit{kruh\textsubscript{2}(-)bʰ-ó}- basic to a *\textit{kruh\textsubscript{2}bʰ-i̯e/o}- ‘render covered, heap (upon)’, whose reflex *\textit{krubʰ-i̯e/o}- was the preform of κρύπτω and the model for the *\textit{kreu̯bʰ}- ‘heap’ ancestral to the Tocharian root. On the basis of the causative periphrasis *\textit{kréh\textsubscript{2}-u-h\textsubscript{1} dʰeh\textsubscript{1}-} was formed a *\textit{kr̥h\textsubscript{2}-u-dʰh\textsubscript{1}-ó}- {>}{>} *\textit{kruh\textsubscript{2}(-)dʰh\textsubscript{1}-ó}-, whose resultative patientive interpretation ‘covered’ paved the way for the quasi-Property Conceptual meaning ‘equipped, adorned’. When this adjective formed an s{}-stem substantive *\textit{kréu̯d\textsubscript{2}(-)dʰh\textsubscript{1}-o/es}- ‘equipment, adornment’ > *\textit{kréu̯dʰh\textsubscript{1}-o/es}-, this quasi-PC meaning and the frequency of \textit{s}-stems beside simple thematic presents, both prototypically radical \textit{é}-grade categories, contributed to the formation of a causative *\textit{kréu̯dʰh\textsubscript{1}-e/o}- ‘render equipped, adorned’ derived by suffixal substitution from a new effectively adjectival root.}


\IfFileExists{../localcommands.tex}{
   \addbibresource{../localbibliography.bib}
   \usepackage{tabularx,multicol}
\usepackage{url}
\urlstyle{same}

 
\usepackage{langsci-optional}
\usepackage{langsci-lgr}
\usepackage{langsci-gb4e}
\usepackage[linguistics,edges]{forest}
\usepackage{tikz-qtree}
\usetikzlibrary{positioning}
\usetikzlibrary{patterns.meta}
\usepackage{multirow}
\usepackage{pgfplots}
\usepackage{setspace}
\usepackage{amsmath,stackengine}
% \usepackage{fontspec}
% \usepackage{tipa} % TIPA is banned
\usepackage{langsci-textipa}
\usepackage{langsci-branding}
\usepackage{longtable}
\usepackage{tabto}

\usepackage{enumitem}

%\usepackage{lscape} %Querformat chapter 4

   % \newcommand*{\orcid}{}


\makeatletter
\let\thetitle\@title
\let\theauthor\@author
\makeatother

\newcommand{\togglepaper}[1][0]{
   \bibliography{../localbibliography}
   \papernote{\scriptsize\normalfont
     \theauthor.
     \titleTemp.
     To appear in:
    Laura Grestenberger, Viktoria Reiter \& Melanie Malzahn (eds.).
    \emph{Deadjectival verb formation in Indo-European and beyond}.
     Berlin: Language Science Press. [preliminary page numbering]
   }
   \pagenumbering{roman}
   \setcounter{chapter}{#1}
   \addtocounter{chapter}{-1}
}

\newbool{bookcompile}
\booltrue{bookcompile}
\newcommand{\bookorchapter}[2]{\ifbool{bookcompile}{#1}{#2}}

%Sonderzeichen
%idg palatales k
\newcommand{\kinvbreve}{k\raisebox{0.6ex}{\hspace{-0.05em}\char"0311}}

\newcommand{\jac}{%
  \textit{\char"0237}% italic dotless j (ȷ)
  \hspace{-0.01em}% fine-tune horizontal spacing
  \raisebox{0.01ex}{\char"0301}% raised combining acute
}

%Baltic circumflex l
\newcommand\ltilded{%
\hspace{-0.2em}%
  \stackon[-1.5pt]{l}{\smash{\kern0.2em\~{}}}%
  \hspace{-0.2em}%
}

%Slav. wedge+ double acute
\newcommand{\HighH}[1]{%
    \stackon[-1.2ex]{#1}{\kern0.1em\H{}\kern-0.1em}%
}

%Lith dot tilde VICKY diese b.-sl. Sonderzeichen BRINGEN MICH UM!!
\newcommand{\tildedot}[1]{%
    \ensurestackMath{%
        \stackon[-0.27ex]{% Tilde layer 
            \stackon[-0.15ex]{% Dot layer 
                #1%
            }{\hspace{0.1em}\cdot\hspace{-0.13em}}%
        }{\hspace{0.1em}\text{\~{}}\hspace{-0.13em}%
    }%
}}

%Lith dot acute
\newcommand{\dotacute}[1]{%
    \ensurestackMath{%
        \stackon[-1.1ex]{% acute layer 
            \stackon[-0.15ex]{% Dot layer 
                #1%
            }{\hspace{0.1em}\cdot\hspace{-0.12em}}%
        }{\hspace{0.1em}\text{\'{}}\hspace{-0.12em}%
    }%
}}



% \setlength{\emergencystretch}{2pt} 

 \pgfplotsset{compat=1.18}


\defcitealias{AbB_1}{AbB 1}
\defcitealias{AbB_7}{AbB 7}
\defcitealias{ABL_1}{ABL 1}
\defcitealias{Vries1977}{AEW}
\defcitealias{ALEW}{ALEW}
\defcitealias{AnGr}{AnGr}
\defcitealias{ARM_1}{ARM 1}
\defcitealias{BT}{BT}
\defcitealias{black_concise_2000}{CDA}
\defcitealias{CHD}{CHD}
\defcitealias{CT28}{CT 28}
\defcitealias{CT34}{CT 34}
\defcitealias{CT39}{CT 39}
\defcitealias{Cleasby1874}{CV}
\defcitealias{Beekes2010}{EDG}
\defcitealias{eDiAna}{eDiAna}
\defcitealias{eDIL}{eDIL}
\defcitealias{EDL}{EDL}
\defcitealias{Kroonen2013}{EDPG}
\defcitealias{ESJS}{ESJS}
\defcitealias{EVP}{EVP}
\defcitealias{EWA1988}{EWA}
\defcitealias{EWAI}{EWAia I}
\defcitealias{EWAII}{EWAia II}
\defcitealias{EWGP}{EWGP}
\defcitealias{GDLI1961}{GDLI}
\defcitealias{GPC}{GPC}
\defcitealias{GRADIT1999}{GRADIT}
\defcitealias{HW}{HW}
\defcitealias{IEW}{IEW}
\defcitealias{KAH_2}{KAH 2}
\defcitealias{KAR_1}{KAR 1}
\defcitealias{KAR_2}{KAR 2}
\defcitealias{KEWA}{KEWA}
\defcitealias{Labat_TDP}{Labat TDP}
\defcitealias{lambert_BWL}{Lambert BWL}
\defcitealias{LEIA}{LEIA}
\defcitealias{Rix2001}{LIV\textsuperscript{2}}
\defcitealias{LIVaddenda}{LIV\textsuperscript{2} {A}ddenda}
\defcitealias{LKG1971}{LKG}
\defcitealias{LKZ}{LKŽ}
\defcitealias{LP}{LP}
\defcitealias{MhdGr}{MhdGr}
\defcitealias{MLLVG1959}{MLLVG}
\defcitealias{NIL}{NIL}
\defcitealias{OgnS}{OgnS}
\defcitealias{ONP}{ONP}
\defcitealias{RIMA_2}{RIMA 2}
\defcitealias{TIM_2}{TIM 2}
\defcitealias{WAP}{WAP}
\defcitealias{YOS_2}{YOS 2}






\newindex{wan}{wdx}{wnd}{Form index}
\newcommand{\iw}[2]{\index[wan]{#1!{#2}\xspace}}
\newcommand{\iwi}[3][]{\index[wan]{#2!{#3}\xspace#1}{#3}}
\newcommand{\iwit}[3][]{\rule{0pt}{0pt}#1\index[wan]{#2!{#3}}{#3}}


\newindex{van}{vdx}{vnd}{Form index}
\newindex{ban}{bdx}{bnd}{Book index}

% \newcommand{\iw}[3][]{\index[wan]{#2!\textit{#3}}}


% \providecommand{\iwi}[3][]{\iw{#3}\textit{#3}}


\newcommand{\indexnote}[1]{{\footnotesize\sffamily\raggedright\normalfont\color{gray}#1}}

   %% hyphenation points for line breaks
%% Normally, automatic hyphenation in LaTeX is very good
%% If a word is mis-hyphenated, add it to this file
%%
%% add information to TeX file before \begin{document} with:
%% %% hyphenation points for line breaks
%% Normally, automatic hyphenation in LaTeX is very good
%% If a word is mis-hyphenated, add it to this file
%%
%% add information to TeX file before \begin{document} with:
%% %% hyphenation points for line breaks
%% Normally, automatic hyphenation in LaTeX is very good
%% If a word is mis-hyphenated, add it to this file
%%
%% add information to TeX file before \begin{document} with:
%% \include{localhyphenation}
\hyphenation{
    par-a-digm non-de-adj-ect-iv-al -pre-sent -pre-sents for-ma-tion-al in-do-ger-ma-ni-sche in-do-ger-ma-ni-schen Rei-chert Rei-ter dia-chrony ety-mo-lo-gi-sches Ost-ger-ma-ni-schen ger-ma-ni-schen alt-in-di-schen mor-pho-lo-gi-scher Alt-indo-ari-sch-en lexico-graphy Schrij-ver
}
% aktuell in Bibliographie: stud-ies, histor-ical Soci-età
\hyphenation{
    par-a-digm non-de-adj-ect-iv-al -pre-sent -pre-sents for-ma-tion-al in-do-ger-ma-ni-sche in-do-ger-ma-ni-schen Rei-chert Rei-ter dia-chrony ety-mo-lo-gi-sches Ost-ger-ma-ni-schen ger-ma-ni-schen alt-in-di-schen mor-pho-lo-gi-scher Alt-indo-ari-sch-en lexico-graphy Schrij-ver
}
% aktuell in Bibliographie: stud-ies, histor-ical Soci-età
\hyphenation{
    par-a-digm non-de-adj-ect-iv-al -pre-sent -pre-sents for-ma-tion-al in-do-ger-ma-ni-sche in-do-ger-ma-ni-schen Rei-chert Rei-ter dia-chrony ety-mo-lo-gi-sches Ost-ger-ma-ni-schen ger-ma-ni-schen alt-in-di-schen mor-pho-lo-gi-scher Alt-indo-ari-sch-en lexico-graphy Schrij-ver
}
% aktuell in Bibliographie: stud-ies, histor-ical Soci-età
   \boolfalse{bookcompile}
   \togglepaper[7]%%chapternumber
}{}

\begin{document}
\SetupAffiliations{mark style=none}
\maketitle





\section{Introduction: from specific to general}

The purpose of this contribution is, on the one hand, to apply a set of theories to explain precisely how the members of a specific set of Indo-European data are historically related and, on the other, to highlight how such explanations offer theoretical implications for our general understanding of various processes whereby verbal stems are formed on the basis of \isi{adnominal}s,\footnote{I use “\isi{adnominal}” in a general sense to refer (with the exception of finite verbs) to any element attributed to or predicated of a substantive.} whether nominal stems or semantically adjectival\footnote{For the derivational behavior of such roots, see \citealt{Nussbaum2022}.} roots. It will be argued that this crucial distinction between the two types of de-adjectival verbs, namely, stem-derived\is{derivation!stem-based} versus root-derived,\is{derivation!deradical} is reflected in the prehistory of, most notably, Greek \iwi{AG}{κρύπτω} ‘cover, hide, bury’ and Old English \iwi{OE}{\textit{hréodan}*} ‘deck, adorn’, which have been assumed to be related in terms of root-etymology (\citetalias[ 616--617]{IEW}), but whose exact derivational prehistory is unclear. The goal of this paper is to clarify that history by (1) identifying the primitive root as \iwi{PIE}{*\textit{kreh\textsubscript{2}}-} (\citetalias[ 367]{Rix2001}), which seems to have meant ‘cover, heap’; (2) proposing that a \textit{u}{}-stem nominal derivative of that root was used in the predicative\is{predicative!instrumental} instrumental \isi{periphrasis} with the light verbs \iwi{PIE}{*\textit{dʰeh\textsubscript{1}}-} ‘put, make’ (\isi{causative}) and \iwi{PIE}{*\textit{bʰuh\textsubscript{x}}-} ‘become’ (\isi{anticausative}); and (3)  arguing for the derivation, on basis of the respective periphrases,\is{periphrasis} of adjectival nominal compounds\is{compounding} whose own derivatives illustrate the interplay of morphology and lexical semantics in de-adjectival derivation. 

Our set of data, derived largely from the material under Pokorny’s \mbox{\iwi{PIE}{*\textit{krā[u]}-} :} \textit{krəu}- : \textit{krū̆-} ‘aufeinander, auf einen Haufen legen, zudecken, verbergen’ (\citetalias[ 616--617]{IEW}), includes Greek \iwi{AG}{κρύπτω} ‘cover, hide, bury’ (: \iwi{AG}{κρύφα} ‘secretly’ : \iwi{AG}{κρυφ-}), Tocharian B \iwi{TB}{\textit{krauptär}}~‘gathers’ (TA \iwi{TA}{\textit{krop}-} ‘id.’), OCS \textit{kryjǫ}, \textit{kryti}\iw{OCS}{\textit{kryti, kryjǫ}} ‘cover’, Lithuanian \textit{kráuju}, \textit{kráuti}\iw{Lith}{\textit{kráuti}, 1sg. prs. \textit{kráuju}} ‘pile up’ (Latv.\ \textit{k\c{r}aũju}, \textit{k\c{r}aũt}\iw{Latv}{\textit{k\c{r}aũt, k\c{r}aũju}} ‘id.’), and Old English \iwi{OE}{\textit{hréodan}*} ‘deck, adorn’ (pret. \textit{hréad} : pret. part. \textit{gehroden}, ON \iwi{ON}{\textit{hrjóða}*} ‘id.’ : \iwi{ON}{\textit{hroðinn}} ‘painted’, \iwi{ON}{\textit{gull-roðinn}} ‘gilt’). Given that all the above-listed data bear segments reflecting the sequence \iwi{PIE}{*\textit{kru}-} and convey lexical meanings involving the notion of covering or placement together, it remains appropriate to maintain that Pokorny is basically correct in supposing that these Greek, Tocharian, Slavic, Baltic, and Germanic forms ultimately share a root reflected in a series of roots reconstructible as secondary.\is{root!secondary}  The question that confronts us if we accept this supposition is the origin of the material following the reflexes of \iwi{PIE}{*\textit{kru}-} in each subset.

\section{The existence of *\textit{kreu̯bʰ-}, *\textit{kreu̯h\textsubscript{x}}-, and *\textit{kreu̯dʰ}-}

The first consists of Greek \iwi{AG}{κρύπτω}  and Tocharian B \iwi{TB}{\textit{kraup}-} (: \iwi{TB}{\textit{kraupe}} m.\ ‘group’), the latter of which, since it forms a class II present, a class II subjunctive, and a class I preterite (\citealt{Adams2013}: 236--238), is likely, in light of the oddity of that pattern, to be inherited (cf.\ athematic Tocharian A \iwi{TA}{\textit{kroptär}} class\is{root!present} I pres., 3pl.\ \mbox{\textit{kropäntär},} \citealt[614--616]{Malzahn2010}). What this comparison allows us to infer is, therefore, that \iwi{AG}{κρύπτω} and \iwi{TB}{\textit{krauptär}} reflect derivatives of a secondary root\is{root!secondary} at least as old as Proto-Nuclear Indo-European, namely a \iwi{PIE}{*\textit{kreu̯bʰ}-} ‘heap, strew, cover’ which formed, in descriptive terms, a radical zero-grade \isi{*-\textit{i̯e/o}-present} \iwi{PIE}{*\textit{krubʰ-i̯e/o}-} ancestral to κρύπτω as well as what appears to be a radical \textit{ó}-grade uncharacterized present ancestral to \iwi{TB}{\textit{krauptär}} (< Transponat \iw{PIE}{*\textit{kreu̯bʰ}-} *\textit{króu̯bʰ-e-tor}). Even if we attribute this root to PIE, it is not, in accordance with our hypothesis, a primary product of reconstruction, as one would assume for a lemma in \citetalias{Rix2001}. To deal with our second subset, which includes OCS \textit{kryti}\iw{OCS}{\textit{kryti, kryjǫ}} ‘cover’ and Lithuanian \textit{kráuti}\iw{Lith}{\textit{kráuti}, 1sg. prs. \textit{kráuju}} ‘pile up’, LIV\textsuperscript{2} reconstructs a \iwi{PIE}{*\textit{kreu̯h\textsubscript{x}}-} ‘aufhäufen, bedecken’ (\citetalias[ 371]{Rix2001}), from which were formed, in descriptive terms, the \iwi{PIE}{*\textit{kruh\textsubscript{x}-i̯e/o}-} ancestral to \textit{kryjǫ}\iw{OCS}{\textit{kryti, kryjǫ}} and the \iwi{PIE}{*\textit{kreu̯h\textsubscript{x}-i̯e/o}-} indirectly ancestral to \textit{kráuju}\footnote{On the acute intonation, see \citealt{VillanuevaSvensson2014}.}\iw{Lith}{\textit{kráuti}, 1sg. prs. \textit{kráuju}} (cf.\ dial.\ Lith.\ \iwi{Lith}{\textit{kriáuju}} : Latv.\ \textit{k\c{r}aũju}\iw{Latv}{\textit{k\c{r}aũt, k\c{r}aũju}}).

While LIV\textsuperscript{2} ignores \iwi{AG}{κρυφ-} and \iwi{TB}{\textit{kraup}-} encountered in Pokorny, it does take account of the Germanic data via the assumption that OE \iwi{OE}{\textit{hréodan}*} ‘bedecken’ reflects a \isi{*-\textit{dʰe/o}-present} \iwi{PIE}{*\textit{kréu̯h\textsubscript{x}-dʰe/o}-}. Though we have good reason to believe that there was indeed a morphological boundary before that dental segment at some historical stage, it appears premature to characterize \iwi{OE}{\textit{hréodan}*} as the reflex of a present in \iwi{PIE}{*-\textit{dʰe/o}-}.\is{*-\textit{dʰe/o}-present} Unlike, for example, \iwi{AG}{βαρύθω} ‘get weighed down (intr.)’, which is formed, with the reflex of a descriptive \iwi{PIE}{*-\textit{dʰe/o}-}, from the actual stem of an adjective (\iwi{AG}{βαρύς} ‘heavy’), the Proto-Germanic \iwi{PGmc}{*\textit{χreuðaną}} ‘deck, equip, adorn’ ancestral to \iwi{OE}{\textit{hréodan}*} behaves as a root-based formation\is{derivation!deradical} parallel to other derivatives of its root (: \iwi{PGmc}{*\textit{χrauðō}-} [ON \iwi{ON}{\textit{hrauð}} f.\ ‘chainmail shirt’, OE \iwi{OE}{\textit{earm-hréad}} ‘arm-ornament’] : \iwi{PGmc}{*\textit{χrusti}-} [OHG \iwi{OHG}{\textit{hrust}} ‘armor’, OE \iwi{OE}{\textit{hyrst}} ‘decoration, armor’]). In this formal respect, \iwi{PGmc}{*\textit{χreuðaną}} rather resembles \iwi{AG}{βρῑ́θω} ‘become weighty’, which, though historically related to βαρ- \iw{AG}{βαρύς} ‘heavy’ (v.\ \citealt[94--95]{Rasmussen1989}), is formed from its own root (: βρῑθ-ύς \iw{AG}{βρῑθύς} ‘heavy’, βρῖθ-oς \iw{AG}{βρῖθoς} ‘weight’). While it is possible to claim that the \iwi{PGmc}{*\textit{-ði/a}-} of \iwi{PGmc}{*\textit{χreuðaną}} reflects an erstwhile presential suffix that formed a stem eventually reanalyzed as a (secondary) \iwi{PIE}{*\textit{kreu̯dʰ}-} ‘cover’,\is{root!secondary} such a claim is of no profit in the absence of an attempt to explain the \iwi{PIE}{*\textit{kreu̯dʰ}-} in question, which must have existed at some stage, in light of any explanations of the origins of the \iwi{PIE}{*\textit{kreu̯bʰ}-} and \iwi{PIE}{*\textit{kreu̯h\textsubscript{x}}-} to which it is assuredly related.

\section{Evidence for primary *\textit{kreh\textsubscript{2}-} ‘heap, strew, cover’}
\largerpage
The initial clue serving in an explanation of how the three secondary roots,\is{root!secondary} \iwi{PIE}{*\textit{kreu̯bʰ}-}, \iwi{PIE}{*\textit{kreu̯h\textsubscript{x}}-}, and \iwi{PIE}{*\textit{kreu̯dʰ}-}, are related is in fact contained in Pokorny’s entry. The brackets of his \iwi{PIE}{*\textit{krā[u]}-}, the first-listed variant of the root, betray the inkling that the truly primary form of the root did not include the \iwi{PIE}{*-\textit{u}-}. Especially in light of the Latvian \textit{krãju}, \textit{krât}\iw{Latv}{\textit{krât, krãju}} ‘put aside, accumulate, amass’ listed under \iwi{PIE}{*\textit{krā[u]}-}, it is deeply appealing to reconstruct the truly primary form  as a \iwi{PIE}{*\textit{kreh\textsubscript{2}}-} ‘heap, strew, cover’, which \citetalias[ 367]{Rix2001} quizzically represents as a \iw{PIE}{*\textit{kreh\textsubscript{2}}-} ?*\textit{kreh\textsubscript{2}}- ‘aufhäufen, sammeln’ supposedly reflected solely in Balto-Slavic. According to Pokorny, one may group \textit{krãju}, \textit{krât}\iw{Latv}{\textit{krât, krãju}} together with OCS \textit{kradǫ}, \textit{krasti}\iw{OCS}{\textit{krasti, kradǫ}} ‘steal’ “mit präsensbildendem d” and, with an apparent \textit{p}-extension, Lithuanian \iwi{Lith}{\textit{krópti}} ‘id.’ and Latvian \textit{krâpju}, \textit{krâpu}, \textit{krâpt}\iw{Latv}{\textit{krâpt, krâpju, krâpu}} ‘cheat, deceive’. Though the origin of such formants or extensions remains to be determined, what is clear is the likelihood that these verbs, given their senses ‘steal’ or ‘deceive’, should be related to the one meaning ‘amass’. Since accumulation generally involves the placement or position of a substance or set of objects over and upon an entity, and since the entity positioned over and upon another entity may serve to effect the latter’s non-stimulus status (e.g., soil of a mound covering a dead body and/or treasure), it is likely that the senses of theft and deception proper to \textit{krasti},\iw{OCS}{\textit{krasti, kradǫ}} \iwi{Lith}{\textit{krópti}}, and \textit{krâpt}\iw{Latv}{\textit{krâpt, krâpju, krâpu}} developed via ‘hide, conceal’ from ‘heap, strew, cover’ (cf.\ PDE \textit{bury}; OIr.\ \iwi{OIr}{\textit{arcelim}} ‘steal’ : \iwi{OIr}{\textit{ceilid}} ‘hides’, \citealt[791]{Buck1949}). 

Some evidence for this original sense of accumulation proper to \iwi{PIE}{*\textit{kreh\textsubscript{2}}-} may also be adduced from Greek. Beside evidence of a \iwi{AG}{κλῶμαξ}, -ακoς, m.\ ‘heap of stones, rocky place’ and a \iwi{AG}{κλωμακόεις} ‘rocky’ (B 729), Hesychius attests to a κρώμαξ· σωρὸς λίθων \iw{AG}{κρώμαξ, κρῶμαξ} and a \iwi{AG}{κρωμακόεν}· κρημνῶδες. While Beekes (\citetalias[ 720]{Beekes2010}), rejecting the association of the former set with \iwi{AG}{κλάω} ‘break’, considers the variation between -λ- and -ρ- to be evidence that both sets are of non-Indo-European origin, an evaluable hypothesis would take them back to two PIE roots. Before considering the question of the roots themselves, we should first observe that the sequence -μακ- points to derivation from stems in \iwi{AG}{-μα} (cf.\ \iwi{AG}{ἕρμαξ}, -ακoς f.\ ‘heap of stones, cairn’ $\leftarrow$ \iwi{AG}{ἕρμα} ‘prop, support’ [Hom.\+]). By comparison to Latin \iwi{Latin}{\textit{rūpēs}} f.\ ‘rock’ (: \iwi{Latin}{\textit{rumpō}} ‘break’), it is probable that, in the case of \iwi{AG}{κλῶμαξ}, the μα-stem \iw{AG}{-μα} in question was indeed akin to \iwi{AG}{κλάω}. If, however, the earlier meaning of that μα-stem was ‘heap’, it is not inconceivable that its preform was derived from *\textit{kleh\textsubscript{2}-} ‘heap up, put’ (\iwi{PIE}{*\textit{k\textsuperscript{(u̯)}leh\textsubscript{2}}-} ‘hinbreiten, hinlegen’, \citetalias[ 362]{Rix2001}), evidence of whose verbal derivatives,\is{derivation!verbal} namely Lithuanian \textit{klóju}, \textit{klóti}\iw{Lith}{\textit{klóti}, 1sg. prs. \textit{klóju}} ‘spread, lay down’ (Latv.\ \textit{klâju}, \textit{klât}\iw{Latv}{\textit{klât, klâju}} ‘id.’) and OCS \textit{kladǫ}, \textit{klasti}\iw{OCS}{\textit{klasti, kladǫ}} ‘load, lay’, incidentally resembles that of \iwi{PIE}{*\textit{kreh\textsubscript{2}}-} ‘heap, strew’. In the case of the base of κρώμαξ,\iw{AG}{κρώμαξ, κρῶμαξ} it is fairly certain that its meaning was ‘heap’.

  Two additional data deserve mention: \iwi{AG}{Kρῶμνα}, -ης (· πόλις   \mbox{Παφλαγoνίας} B 855 [Hsch.]; a toponym) and the reportedly Paphlagonian \iwi{AG}{κρωμακωτός} ‘rocky and steep’. Given the identity of place as well as toponomastic practice (cf.\ \iwi{AG}{Αἴπεια} ‘Hightown’ : \iwi{AG}{αἰπύς, αἰπός} ‘high, lofty’), it is likely that the base of the \textit{devī́}-derivative ancestral to \iwi{AG}{Kρῶμνα} was reflected as the μα-stem basic to the κρῶμαξ  \iw{AG}{κρώμαξ, κρῶμαξ}  from which \iwi{AG}{κρωμακωτός} was derived.  In all likelihood, this μα-stem would go back to a \mbox{*\textit{-men}-} stem derived from \iwi{PIE}{*\textit{kreh\textsubscript{2}}-} ‘heap, strew’.

  Comparable in form to the \iwi{PIE}{*\textit{póh\textsubscript{2}-mn̥}} ‘protection’ (: \iwi{PIE}{*\textit{peh\textsubscript{2}}-} ‘cover, protect, care for’) ancestral to \iwi{AG}{πῶμα} ‘lid, cover’ (Hom.\+),\footnote{Analysis of \iwi{AG}{πῶμα} ‘lid, cover’ as the reflex of an *\textit{ó}/\textit{é} \isi{acrostatic} derives from Nussbaum (p.c.).} this *-\textit{men-} stem would have been an *\textit{ó}/\textit{é} \isi{acrostatic}. In addition to the association between κρῶμαξ\iw{AG}{κρώμαξ, κρῶμαξ} and \iwi{PIE}{*\textit{kreh\textsubscript{2}}-} effected by the attribution of the *-\textit{men-} stem to this type, the reconstruction of a \iwi{PIE}{*\textit{króh\textsubscript{2}-mn̥}/\textit{kréh\textsubscript{2}-mn̥}-} ‘heaping, heap’ has the potential advantage of providing a way to account for the origin of \iwi{AG}{κρημνός} ‘overhanging bank, cliff’. Since the appearance of its -η- in Pindar appears to preclude an *\textit{-ā-} in the preform (\citetalias[ 777]{Beekes2010}), it is conceivable that \iwi{AG}{κρημνός} goes back to a \isi{genitival} \iwi{PIE}{*\textit{krēh\textsubscript{2}-mn-ó}-} ‘of the heaping, heaped’ derived from \iw{PIE}{*\textit{króh\textsubscript{2}-mn̥}/\textit{kréh\textsubscript{2}-mn̥}-} \mbox{*\textit{króh\textsubscript{2}-mn̥}/}\textit{kréh\textsubscript{2}-mn̥}- ‘heaping, heap’. Though this scenario has some problems,\footnote{One would have to assume that this derivation took place at a very early date in the protolanguage and that one of the nasals failed to be deleted as expected or was restored.} it does provide \iwi{AG}{κρημνός} with a semantically straightforward etymology (cf.\ Hitt. \iwi{Hitt}{\textit{wappu}-} ‘riverbank’ < \mbox{*\textit{(h\textsubscript{2})u̯óp-u-}/}\textit{h\textsubscript{2}u̯ép-u-} \iw{PIE}{*\textit{(h\textsubscript{2})u̯ó/ép-u}-} \mbox{‘heap’ :} \iwi{PIE}{*\textit{h\textsubscript{2}u̯ep}-} ‘cast, strew’ [Hitt. \iwi{Hitt}{\textit{ḫuwapp-}/\textit{ḫupp}-} ‘cast, throw’ : Ved.\ \iwi{Sanskrit}{\textit{vápati}} ‘strews’], \citealt[125]{Catsanicos1985}; \citealt[176]{Melchert2013}) whose formal features are mostly in line with a simple thematic item conceivably derived from a *-\textit{men-} stem to \iwi{PIE}{*\textit{kreh\textsubscript{2}}-} ‘heap, strew’.

By comparison to the form of Latin \iwi{Latin}{\textit{fās}} ‘morally right’ (< \iwi{PIE}{*\textit{bʰéh\textsubscript{2}-o/es}-} ‘saying, utterance’, \citetalias[ 203]{EDL}), we may hypothesize that \iwi{PIE}{*\textit{kreh\textsubscript{2}}-} also formed an \textit{s}-stem substantive \iwi{PIE}{*\textit{kréh\textsubscript{2}-s}-} ‘heaping, strewing, spreading, covering’ for which there is some indirect evidence in Baltic. Beside Lithuanian \iwi{Lith}{\textit{krū́snis}, -\textit{ies} f.} ‘pile of stones in a field’ (v.\ \citealt[616]{Smoczynski2018}), identifiable as the reflex of a \iwi{PIE}{*\textit{kruh\textsubscript{2}-s-n-i}-} derived from a \isi{possessive} in \iwi{PIE}{*\textit{-nó}-} to an \textit{s}-stem \iwi{PIE}{*\textit{kréu̯h\textsubscript{2}-o/es}-} ‘heap, covering’ to be discussed below (: \iwi{PIE}{*\textit{króu̯h\textsubscript{2}-o}-} ‘covering’ [> OCS \iwi{OCS}{\textit{krovъ}} ‘roof’]), Lithuanian and Latvian respectively attest to a \iwi{Lith}{\textit{krósnis}, -\textit{ies} f.}\ ‘(stone-)stove; pile of stones’ and \iwi{Latv}{\textit{krâsns}} ‘stove’, which \citet[613]{Smoczynski2018} and \citet[419]{Karulis1992} both connect with Latvian \textit{krãju}, \textit{krât}\iw{Latv}{\textit{krât, krãju}} ‘amass’. While Smoczyński’s connection with \textit{krât}\iw{Latv}{\textit{krât, krãju}} and thus \iwi{PIE}{*\textit{kreh\textsubscript{2}}-} is entirely justified, what this scholar analyzes as a suffix \textit{-sni-} is in need of further clarification. If we compare the dialectal Latvian \iwi{Latv}{\textit{krāsts}} m.\ ‘oven’ (: \iwi{Latv}{\textit{krāstît}} ‘collect, hoard’), we may suppose a formation similar to Latin \iwi{Latin}{\textit{fāstus}} ‘right’ ($\leftarrow$ \iwi{Latin}{\textit{fās}} ‘id.’), that is, a \isi{possessive} in *-\textit{tó}-\iw{PIE}{*-\textit{to}-} derived from an \textit{s}-stem basic to a *\textit{-nó-} stem from which the ancestor of this -\textit{sni-} stem was formed. Since that \textit{s}-stem had the sense ‘heaping, heap’, its derivative, namely \iwi{PIE}{*\textit{kr̥h\textsubscript{2}-s-nó}-} (cf.\ \iwi{AG}{λύχνoς} ‘lamp’ < \iwi{PIE}{*\textit{lúk-s-no}-} ‘the lighted’ $\leftarrow$ \iwi{PIE}{*\textit{luk-s-nó}-} ‘lighted’ $\leftarrow$ \iwi{PIE}{*\textit{léu̯k-o/es}-} ‘light’), would have borne a meaning ‘heaped’ subsequently determined by means of the \isi{acrostatic} \textit{i}-stem derivative \iwi{PIE}{*\textit{kréh\textsubscript{2}-s-n-i}-} ‘the heaped, heap’ ancestral to the Baltic cognates. Since the above nominal and verbal evidence supports the reconstruction of \iwi{PIE}{*\textit{kreh\textsubscript{2}}-} as a primary root, the immediate question is the nature of the \iwi{PIE}{*-\textit{u}-} in the secondary roots.\is{root!secondary} As the old extensions invoked by Pokorny have increasingly ceded to the discipline’s growing powers of morphosemantic\is{morphosemantics} analysis, we can assume that the \iwi{PIE}{*-\textit{u}-} is of suffixal origin.

\section{The nature and derivational role of *\textit{kró/éh\textsubscript{2}-u-} ‘heap, covering’}

\subsection{Background}

This assumption opens up two paths of hypothesis: (1) it is a verbal stem formant, in which case it would be the suffix of a \textit{u}-present,\is{*-\textit{u}-present} or (2) it is the suffix of a nominal stem.\ The preferability of the latter hypothesis becomes clear in light of parallelism with the derivational history assigned in previous research to \iwi{AG}{καλύπτω} ‘cover’ \citep{Merritt2021}. This particular history derives its existence from the understanding, originating in and enriched by, respectively, comparative-historical (\citealt[118--126]{Jasanoff1978}; \citeyear{Jasanoff2004}; \citealt{Schindler1980}; \citealt{Balles2006}; \citealt{Rau2009}) and theoretical-typological work (\citealt{Dixon1982,Dixon2004}; \citealt{FrancezKoontz2015}), that PIE was a language in which certain \isi{stative} concepts, such as Property Conceptual (PC)\is{property concept (PC)} states, were disposed to realization via substantives in the instrumental case. In prelinguistic terms, a \isi{stative} concept could be accessed by conceiving of a substance, whether concrete or abstract, as an entity possessed by the being of which the state is predicated. While much attention has been paid to the role of adjectival (i.e., PC)\is{property concept (PC)} abstracts in this means of \isi{stative} predication, we should not forget that the specific \isi{stative}/\isi{anticausative}/\isi{causative}\footnote{I use such terms in a general sense (cf.\ introduction to this volume: 14--16). Accordingly, a \isi{causative} and \isi{anticausative} derived from an adjective or adjectival root\is{root!adjectival} are specifically \isi{factitive} and \isi{fientive} respectively.} construction which inaugurated much of this research tradition, namely Vedic \textit{gúhā as-}/\mbox{\iwi{Sanskrit}{\textit{bhū}-}/}\iwi{Sanskrit}{\textit{dhā}-} \iw{Sanskrit}{\textit{gúhā}} ‘be/become/make hidden’, includes the reflex of a substantive derived ultimately from a root that basically meant ‘cover’ (Lith.\ \iwi{Lith}{\textit{gū̃žti}} ‘[of hens] sit, squat, keep warm’), corresponding to an Activity \is{activity} Concept, rather than a Property Concept.\is{property concept (PC)} Though \iwi{Sanskrit}{\textit{gúhā}} in the sense ‘hidden, secret’ does express (approximately) a PC\is{property concept (PC)} state, its preform, in light of the \isi{metonymy} \textsc{covered for hidden} underlying that sense, must have resembled a verbal \isi{resultative} state, namely ‘covered’ (lit.\ ‘with covering’), predicated of a patient/theme. In like manner, there is, at the heart of the history of \iwi{AG}{καλύπτω}, an \isi{acrostatic} \textit{u}-stem substantive \iwi{PIE}{*\textit{\kinvbreve ó/él-u{}}-} ‘covering’, whose instrumental was used to predicate just such an Activity \is{activity} Concept state. Correspondingly, it seems feasible to suppose that our \iwi{PIE}{*\textit{kreh\textsubscript{2}}-} was ultimately the base of a \iwi{PIE}{*\textit{kró/éh\textsubscript{2}-u}-} ‘heap, covering’. 

\subsection{*\textit{kreh\textsubscript{2}-u-i̯e/o}- ‘provide with a covering’}

In accordance with this hypothesis, it is conceivable that the instrumental of \iwi{PIE}{*\textit{kró/éh\textsubscript{2}-u}-}, namely \iwi{PIE}{*\textit{kréh\textsubscript{2}-u-h\textsubscript{1}}}, was used as an \isi{adnominal} in nominal sentences or (anti)\isi{causative} \is{anticausative}  periphrases\is{periphrasis} with forms of \iwi{PIE}{*\textit{bʰuh\textsubscript{x}}-} ‘become’ and \iwi{PIE}{*\textit{dʰeh\textsubscript{1}}-} ‘put, make’. While the instrumental form would participate in \isi{periphrasis}, it is also conceivable that the meaning of the instrumental could be expressed without the morphology of that case by means, for instance, of the suffix \iwi{PIE}{*-\textit{i̯e/o}-}.\is{*-\textit{i̯e/o}-present} As a denominal verbal \is{denominal!verb} stem formant, \iwi{PIE}{*-\textit{i̯e/o}-} lets the constructional and lexical semantic idiosyncrasies of the base shape the meaning of what it derives. In principle, it does not strictly specify the quantity or quality of the arguments. In practice, however, there seems to be a tendency for transitivity\is{transitive} and for the agentivity \is{agentive} of the argument serving as subject. For example, the derivational semantic history of Latin \iwi{Latin}{\textit{arguō}} ‘demonstrate’ is discernible by supposing that it reflects a \isi{transitive} \isi{*-\textit{i̯e/o}-present} \iwi{PIE}{*\textit{h\textsubscript{2}erg̑-u-i̯e/o}-}\footnote{The cognation of Hittite \textit{arkuwae-\textsuperscript{zi}} \iw{Hitt}{\textit{arkuwae}-} ‘make a plea’, which \citealt{Melchert1998} associated with \iwi{Latin}{\textit{arguō}} (< *\textit{argu-yé/ó-}), would be difficult to maintain, unless one assumes on the basis of \citealt{Kloekhorst2006} and \citeyear[206]{Kloekhorst2008} that, since the root \iwi{PIE}{*\textit{h\textsubscript{2}erg̑}-} of the form ancestral to \textit{arkuwae-\textsuperscript{zi}} \iw{Hitt}{\textit{arkuwae}-} was in \textit{o}-grade (rather than the \textit{e}-grade found elsewhere), the initial \textit{h\textsubscript{2}} was lost.} ‘[\textit{x}] provide [[\textit{y}] [with clarity]]’ derived from \iwi{PIE}{*\textit{h\textsubscript{2}ó/érg̑-u}-} ‘clarity’ in the sense corresponding to its instrumental \iwi{PIE}{*\textit{h\textsubscript{2}érg̑-u\nobreakdash-h\textsubscript{1}}} ‘with clarity’, from which in theory was formed a \textit{t}{}-stem (abstract)\footnote{Cf.\ Hitt. \iwi{Hitt}{\textit{ḫappina}-} ‘rich’ $\rightarrow$ abstract \iwi{Hitt}{\textit{ḫappinatt}-} (c.) ‘wealth’, \citealt{Nussbaum2004}.} substantive basic to the \isi{possessive} adjective ancestral to the “presuffixally lengthened” \isi{perfect} passive \isi{participle} \iwi{Latin}{\textit{argūtus}} ‘demonstrated’ (< Transponat \iw{PIE}{*\textit{h\textsubscript{2}erg̑}-} *\textit{h\textsubscript{2}erg̑-u-h\textsubscript{1}-t-o}-; \citealt{Nussbaum1996}; cf.\ \citetalias[ 53]{EDL}). Given the frugality of this analysis, it is appealing to suppose that \iwi{PIE}{*\textit{kró/éh\textsubscript{2}-u}-} ‘heap, a covering’ was the base of an instrumentative PIE \isi{*-\textit{i̯e/o}-present} \iw{PIE}{*-\textit{i̯e/o}-} \iwi{PIE}{*\textit{kreh\textsubscript{2}-u-i̯e/o}-} ‘provide with a covering’.

\subsection{The genesis of *\textit{kreu̯h\textsubscript{2}}- ‘heap, cover’}


Since this verb’s meaning was equivalent to that of the verb originally derived from \iwi{PIE}{*\textit{kreh\textsubscript{2}}-} ‘heap, strew, cover’, it is conceivable that it was reanalyzed as a primary present\is{derivation!deradical}  derived from a “long-diphthongal” \iwi{PIE}{*\textit{kreh\textsubscript{2}u}-} ‘heap, strew, cover’. Accordingly, the root of this present would have appeared eligible as a base to a *-\textit{tó}- \iw{PIE}{*-\textit{to}-} \isi{verbal adjective}, namely a \iwi{PIE}{*\textit{kr̥h\textsubscript{2}u-tó}-} ‘heaped, covered’. Though not of the canonical structure reflected in Vedic \iwi{Sanskrit}{\textit{pītá}-} ‘drunk’ (< \iwi{PIE}{*\textit{pih\textsubscript{3}-tó}-} < \mbox{\iwi{PIE}{*\textit{ph\textsubscript{3}(-)i-tó}-}),} this \iwi{PIE}{*\textit{kr̥h\textsubscript{2}u-tó}-} may be assumed to have undergone laryngeal metathesis in the same way as the preform of Latin ($\leftarrow$ Sabellic) \iwi{Latin}{\textit{brūtus}} ‘heavy, stupid’, which via \iwi{PIE}{*\textit{g\textsuperscript{u̯}ruh\textsubscript{2}tó}-} would reflect, whatever its exact analysis, a *\textit{\mbox{g\textsuperscript{u̯}r̥h\textsubscript{2}-u-t(-)}ó-} ‘heavy’ related to the preform of \iwi{AG}{βαρύς} ‘id.’. The resulting \iwi{PIE}{*\textit{kruh\textsubscript{2}-tó}-} ‘heaped, covered’, arguably basic to the preform of Lithuanian \iwi{Lith}{\textit{krū́tis} f.}\ ‘pile of stones, stone-oven’ (< \iwi{PIE}{*\textit{kruh\textsubscript{2}-t-i}-}) and directly ancestral to the past passive \isi{participle} of OCS \textit{kryti}\iw{OCS}{\textit{kryti, kryjǫ}} ‘cover’ (cf.\ Russian \iwi{Russ}{\textit{krýtyj}} ‘covered’), appears to have been pivotal in the creation of what we may now rewrite as our secondary\is{root!secondary} \iwi{PIE}{*\textit{kreu̯h\textsubscript{2}}-} ‘heap, cover’ ($\rightarrow$ τόμoς-substantive \iwi{PIE}{*\textit{króu̯h\textsubscript{2}-o}-} ‘covering’ [> OCS \iwi{OCS}{\textit{krovъ}} ‘roof’] : \textit{s}-stem \iwi{PIE}{*\textit{kréu̯h\textsubscript{2}-o/es}-} ‘heap, covering’ $\rightarrow$ \iwi{PIE}{*\textit{kruh\textsubscript{2}-s-ó}-} ‘heaped, cover(ed)’ [$\rightarrow$ \iwi{PIE}{*\textit{króu̯(h\textsubscript{2})-s-o}-} $\rightarrow$ \iwi{PIE}{*\textit{króu̯(h\textsubscript{2})-s-ii̯o}-} > ON \iwi{ON}{\textit{hreysi}} ‘stone heap’] $\rightarrow$ \iwi{PIE}{*\textit{kruh\textsubscript{2}-s-ii̯o}-} $\rightarrow$ \iwi{PIE}{*\textit{krúh\textsubscript{2}-s-ii̯e-h\textsubscript{2}}-} > Russian \iwi{Russ}{\textit{krýša}} ‘roof’). While Slavic points to a radical zero-grade \iwi{PIE}{*\textit{kruh\textsubscript{2}-i̯e/o}-} derived from \iwi{PIE}{*\textit{kreu̯h\textsubscript{2}}-} on the model of \iwi{PIE}{*\textit{kruh\textsubscript{2}-tó}-}, Baltic inherited a radical \textit{e}-grade \isi{*-\textit{i̯e/o}-present} \iwi{PIE}{*\textit{kreu̯h\textsubscript{2}-i̯e/o}-} ‘heap’ ({>}{>} dial.\ Lith.\ \iwi{Lith}{\textit{kriáuju}} : Latv.\ \textit{k\c{r}aũju}\iw{Latv}{\textit{k\c{r}aũt, k\c{r}aũju}}) that arose conceivably by analogy to families involving (as usual) radical zero-grade *-\textit{tó}- \iw{PIE}{*-\textit{to}-} verbal adjectives\is{verbal adjective} beside radical \textit{e}-grade \iwi{PIE}{*-\textit{i̯e/o}-}presents \is{*-\textit{i̯e/o}-present} (e.g., \iwi{PIE}{*\textit{ten-i̯e/o}-} ‘stretch’ [\iwi{AG}{τείνω}, Alb.\ \iwi{Albanian}{\textit{n-den}} ‘stretches out’,  \citetalias[ 627]{Rix2001}] : \iwi{PIE}{*\textit{tn̥-tó}-} ‘stretched’ [Ved.\ \iwi{Sanskrit}{\textit{tatá}-} ‘id.’, \iwi{AG}{τατός} ‘extensible’]).

\section{The origin of κρύπτω}
\subsection{Background}

To understand the emergence of \iwi{PIE}{*\textit{kreu̯bʰ}-} beside the \iwi{PIE}{*\textit{kreu̯h\textsubscript{2}}-} whose origin we have just considered, we should return to the original instrumentative \iwi{PIE}{*\textit{kreh\textsubscript{2}-u-i̯e/o}-} ‘provide with a covering’. Since, by its very nature, this denominal \is{denominal!verb} was a creature of the present system, the non-present system had to be served by periphrases\is{periphrasis} consisting of the instrumental of the denominal’s base and the light verbs \iwi{PIE}{*\textit{dʰeh\textsubscript{1}}-} and \iwi{PIE}{*\textit{bʰuh\textsubscript{x}}-}. Given the necessity of \isi{periphrasis} at some level, it is conceivable that the option to form compounds\is{compounding} based on that \isi{periphrasis} was realized. We may thus suppose that the \isi{anticausative} phrase \iwi{PIE}{*\textit{kréh\textsubscript{2}-u-h\textsubscript{1}}} \textit{bʰuh\textsubscript{x}}- \iw{PIE}{*\textit{bʰuh\textsubscript{x}}-} ‘become covered’ inspired a \isi{resultative} compound adjective\is{compounding} \iwi{PIE}{*\textit{kr̥h\textsubscript{2}-u-bʰuh\textsubscript{x}-ó}-} ‘(having) become heaped, covered’, whose νεoγνός-reflex was a \iwi{PIE}{*\textit{kr̥h\textsubscript{2}-u-bʰ(u̯)-ó}-}. This form was ancestral to an eventually simple thematic adjective \iwi{PIE}{*\textit{kruh\textsubscript{2}bʰ-ó}-}\footnote{For the PIE status of the development *\textit{-bʰ}\textit{u̯-} to *\textit{-bʰ}\textit{-}, see \citealt[312--313]{Merritt2021}.} ‘heaped, covered’ via laryngeal metathesis and morphological \isi{reanalysis}. As a thematic adjective, this \iwi{PIE}{*\textit{kruh\textsubscript{2}bʰ-ó}-} would in theory have been eligible as a base to a \textit{nēwaḫḫi}-type\is{\textit{nēwaḫḫi}-type} \isi{factitive}. What we find instead is the reflex of a \isi{*-\textit{i̯e/o}-present} which, like \iwi{AG}{χαλέπτω} ‘treat harshly’ beside \iwi{AG}{χαλεπός} ‘harsh’, was derived by substitution of *-\textit{i̯e/o}- for the thematic vowel of the \isi{adnominal} base. The derivative would thus have been \iwi{PIE}{*\textit{kruh\textsubscript{2}bʰ-i̯e/o}-} ‘render covered, heap (upon)’.

\subsection{The parallel of πλάσσω ‘form, mould’}

The morphological and phonological nature and fate of this denominal \is{denominal!verb} may be compared to the history of \iwi{AG}{πλάσσω} ‘form, mould’, which, given the sequence πλαθ{}- (: \iwi{AG}{κoρoπλάθoς} ‘figure-modeler’), Pokorny (\citetalias[ 805--807]{IEW}) characterizes as a \isi{*-\textit{i̯e/o}-present} formed from a base containing the remnant of a \textit{dʰ}-formant of a present stem derived ultimately from what would now be represented as \iwi{PIE}{*\textit{pleh\textsubscript{2}}-} ‘broad, flat’ (Hitt. \iwi{Hitt}{\textit{palḫi-}/\textit{palḫai}-} ‘wide, broad’, Lat.\ \iwi{Latin}{\textit{plānus}} ‘flat’, OE \iwi{OE}{\textit{flōr}} f.\ ‘floor’). Since, however, \iwi{AG}{πλάσσω} does reflect a present in *-\textit{i̯e/o}- \is{*-\textit{i̯e/o}-present} (PGk.\ \mbox{\iwi{PG}{*\textit{plətʰ-i̯e/o}-}),} it is preferable to analyze the preform of this verb as denominal,\is{denominal!verb} specifically de-\isi{adnominal} from a *-\textit{dʰh\textsubscript{1}-o-} \iw{PIE}{*-\textit{dʰh\textsubscript{1}-ó}-} compound.\is{compounding} In maintenance of the connection with the descriptive \iwi{PIE}{*\textit{pleh\textsubscript{2}}-} ‘broad, flat’, such an analysis would begin with the derivation of what may be provisionally represented as a root noun\is{root!noun} \iwi{PIE}{*\textit{pléh\textsubscript{2}-}/\textit{pl̥h\textsubscript{2}}-´} ‘the flat, flatness’, whose capacity to modify substantives rendered it eligible for use in the familiar predicative\is{predicative!instrumental} instrumental \isi{periphrasis} (i.e., the \textit{cale-faciō}-type). Accordingly, on the basis of a \isi{causative} phrase \iwi{PIE}{*\textit{pl̥h\textsubscript{2}-éh\textsubscript{1}}}  \textit{dʰeh\textsubscript{1}}-\iw{PIE}{*\textit{dʰeh\textsubscript{1}}-} ‘provide with the flat, flatten’ was formed \iw{PIE}{*\textit{pl̥h\textsubscript{2}(-)dʰh\textsubscript{1}-ó}-}  *\textit{pl̥h\textsubscript{2}-dʰh\textsubscript{1}-ó}- ‘flattened, pressed, kneaded, moulded’, apparently basic to the substantive \iwi{PIE}{*\textit{pl̥h\textsubscript{2}(-)dʰh\textsubscript{1}-e-h\textsubscript{2}}-} ‘the moulded, mould’ ancestral, via the παλάμη- or \textit{palma}-rule (\citealt[157--158]{Rix1996}, \citealt[108--109]{Meiser1998}, \citealt{Hofler2017}, \citealt[119]{Weiss2020}, \citealt[193--194]{Hackstein2021}), to the παλάθη\iw{AG}{παλάθη, παλάθαι} ‘cake of preserved fruit’ adduced by Pokorny (‘flacher Fruchtkuchen’; cf.\ παλάθαι· σύκων μαζία [Hsch.] : \iwi{AG}{μάσσω} ‘knead, mould’). 

The sequence παλαθ- brings up the question of the origin of the radical base πλᾰθ- underlying \iwi{AG}{πλάσσω} and apparent in its nominal relatives (e.g., \iwi{AG}{πλάθανoν} ‘baking-mould’, \iwi{AG}{πλαθά} f.\ ‘modelled figure’). Given that πλᾰθ- does not display a laryngeal reflex, it is necessary to suppose that speakers posited the root on the basis of a form in which the laryngeal was deleted. Since the substitution of \iwi{PIE}{*-\textit{i̯e/o}-} for the thematic vowel of our \iwi{PIE}{*\textit{pl̥h\textsubscript{2}(-)dʰh\textsubscript{1}-ó}-} would have produced a sequence subject to Pinault’s Law and, subsequently, the \textit{Weather} Rule (cf.\ Ion.\ \iwi{AG}{γλᾰ́σση-} ‘tongue’ < *\textit{gl̥h\textsubscript{3}gʰ-i̯eh\textsubscript{2}}- ‘pointed’ [cf.\ \iwi{AG}{γλωχῑ́ς} ‘point, barb’] : *\textit{gl(e)h\textsubscript{3}gʰ-ih\textsubscript{2}}-\iw{PIE}{*\textit{gl(e)h\textsubscript{3}gʰ-ih\textsubscript{2}}-/*\textit{gl̥h\textsubscript{3}gʰ-i̯eh\textsubscript{2}}-} \mbox{[: \iwi{AG}{γλῶττα}} ‘tongue’]; \citealt{Neri2017}; \citealt[123]{Weiss2020}), it is appropriate to suppose that \iwi{PIE}{*\textit{pl̥h\textsubscript{2}(-)dʰh\textsubscript{1}-i̯e/o}-} ‘get x moulded’ was reflected as an eventual \iwi{PIE}{*\textit{pl̥dʰ-i̯e/o}-} directly ancestral to \iwi{AG}{πλάσσω}, whose meaning ‘form’ resulted from a metonymic\is{metonymy} extension on the part of \iwi{PIE}{*\textit{pl̥h\textsubscript{2}(-)dʰh\textsubscript{1}-ó}-} from ‘flattened, pressed’ to ‘kneaded, moulded’ via the notion of applying pressure to change or create a surface. Morphologically, the result is clear: the \isi{*-\textit{i̯e/o}-present} \iwi{PIE}{*\textit{pl̥dʰ-i̯e/o}-} was reanalyzed as primary\is{derivation!deradical}  from a \iwi{PIE}{*\textit{pledʰ}-} ‘form, shape’ that could now form derivatives of its own, such as the τόμoς-type \iwi{PIE}{*\textit{plódʰ-o}-} ‘formation’ arguably ancestral to OCS \iwi{OCS}{\textit{plodъ}} ‘fruit, fetus’ (cf.\ \iwi{PIE}{*\textit{kreu̯h\textsubscript{2}}-} ‘heap, cover’ $\rightarrow$ \iwi{PIE}{*\textit{króu̯h\textsubscript{2}-o}-} ‘covering’ > OCS \iwi{OCS}{\textit{krovъ}} ‘roof’).

\subsection{The origin of *\textit{kreu̯bʰ}-}

If we apply the comparandum that is \iwi{PIE}{*\textit{pl̥h\textsubscript{2}(-)dʰh\textsubscript{1}-i̯e/o}-} to our \iwi{PIE}{*\textit{kruh\textsubscript{2}bʰ-i̯e/o}-} ‘render covered, heap (upon)’, the result is equally clear. Once the \textit{Weather} Rule deleted the laryngeal, the sequence consisted of a suffix \iwi{PIE}{*-\textit{i̯e/o}-} and a base \iwi{PIE}{*\textit{krubʰ}-} interpretable as a root in its own right, namely \iwi{PIE}{*\textit{kreu̯bʰ}-} ‘heap, cover’. At this point, we can revisit the thematic Tocharian B \iwi{TB}{\textit{krauptär}} ‘gathers’ in light of its unambiguously athematic cognate\is{root!present} \iwi{TA}{\textit{kroptär}} in Tocharian A (class I pres., 3 pl.\ \mbox{\textit{kropäntär},} \citealt[614--616]{Malzahn2010}). Given this discrepancy, it appears probable that the athematicity of the Tocharian A form is a retention (Jasanoff p.c.). Given, moreover, the likelihood of this retention and the possibility that \iwi{TB}{\textit{kraup}-} and \iwi{TA}{\textit{krop}-} reflect an \textit{o}-grade, it is probable that Proto-Tocharian inherited the reflex of a \isi{\textit{molō}-present} (\citealt[70--79]{Jasanoff2003}) \iwi{PIE}{*\textit{króu̯bʰ-}/\textit{kréu̯bʰ}-} ‘heap up together’ derivationally parallel to the τόμoς-type \iwi{PIE}{*\textit{króu̯bʰ-o}-} ‘accumulation’ apparently ancestral to TB \iwi{TB}{\textit{kraupe}} ‘group’. In other words, once a \iwi{PIE}{*\textit{kreu̯bʰ}-} ‘heap, cover’ was posited on the basis of a historically secondary formation\is{derivation!stem-based}  synchronically interpretable as primary, namely a \isi{*-\textit{i̯e/o}-present}, the formation of primary stems,\is{derivation!deradical}  namely a \isi{\textit{molō}-present} and τόμoς-type nominal, became possible.\footnote{A similar account may be assigned to the preform of Proto-Germanic \iwi{PGmc}{*\textit{walð-i/a}-} : Proto-Slavic \iwi{PSl}{*\textit{vold-e/o}-} : Proto-Baltic \iwi{PB}{*\textit{véld}-} ‘be powerful, rule’, which \citet[129--135]{RothsteinDowden2022} has analyzed as a \isi{\textit{molō}-present} formed with suffixal *-\textit{dʰ-} (\iwi{PIE}{*\textit{(h\textsubscript{2})u̯ólh\textsubscript{x}-dʰ-}/*\textit{(h\textsubscript{2})u̯élh\textsubscript{x}-dʰ}-} ‘be powerful’).}

\section{The origin of *\textit{χreuðaną}}
\subsection{Introduction}

A similar kind of root-based\is{derivation!deradical} account is crucial to explaining the origin of secondary\is{root!secondary} \iwi{PIE}{*\textit{kreu̯dʰ}-} and Proto-Germanic \iwi{PGmc}{*\textit{χreuðaną}} ‘deck, equip, adorn’. In light of our account of the origin of the secondary\is{root!secondary} \iwi{PIE}{*\textit{kreu̯bʰ}-} in the \isi{anticausative} phrase \iwi{PIE}{*\textit{kréh\textsubscript{2}-u-h\textsubscript{1}}} \textit{bʰuh\textsubscript{x}}-\iw{PIE}{*\textit{bʰuh\textsubscript{x}}-} ‘become covered’, it is reasonable to suppose that \iwi{PIE}{*\textit{kreu̯dʰ}-}, provisionally definable as ‘cover’, originated ultimately in the \isi{causative} variant of the same construction, namely \iwi{PIE}{*\textit{kréh\textsubscript{2}-u-h\textsubscript{1}}} \textit{dʰeh\textsubscript{1}}-\iw{PIE}{*\textit{dʰeh\textsubscript{1}}-} ‘provide with covering’, on the basis of which a compound\is{compounding} adjective \iwi{PIE}{*\textit{kr̥h\textsubscript{2}-u-dʰh\textsubscript{1}-ó}-} was formed. Though this adjective could certainly convey the meaning ‘providing with covering’ as a function of the agentivity \is{agentive} assignable by the \isi{causative} light verb, the necessary existence of a patient/theme in the structure of \iwi{PIE}{*\textit{dʰeh\textsubscript{1}}-} would always render possible the interpretation of \iwi{PIE}{*\textit{kr̥h\textsubscript{2}-u-dʰh\textsubscript{1}-ó}-} as an adjective of patientive and \isi{resultative} sense equivalent to its sister \iwi{PIE}{*\textit{kr̥h\textsubscript{2}-u-bʰuh\textsubscript{x}-ó}-} ‘(having) become heaped, covered’. We may thus suppose that \iwi{PIE}{*\textit{kr̥h\textsubscript{2}-u-dʰh\textsubscript{1}-ó}-} also meant ‘provided with cover, rendered covered’.

\subsection{\textsc{covered for \textit{x}}}

This patientive and \isi{resultative} sense ‘covered’ is a nucleus disposed, via the basic notion of positioning one entity over and upon another entity, to generate by \isi{metonymy} a variety of senses, some of which correspond approximately to Property Concepts.\is{property concept (PC)} For example, the sense of Vedic \iwi{Sanskrit}{\textit{gúhā}} ‘secret’ is, as discussed above, the product of \textsc{covered for hidden}. The structurally and semantically interesting feature of this \isi{metonymy} is its introduction of the end-oriented notion that the action relates to a third entity, which in this case is a non-experiencer, beside the patient/theme and implied agent. In other words, when a verb that means ‘cover’ is used in the sense ‘hide’, there arises the implication that the end or purpose of the agent’s action is to prevent the patient/theme from becoming a stimulus to some additional person (third entity), whose existence is implied, and which may even surface as an participant, as in (agent) \textit{x}’s hiding of \textit{y} (non-stimulus) from \textit{z} (non-experiencer). The \isi{metonymy} \textsc{covered for hidden} that generates the implication of this third entity is structurally comparable to two other metonymies,\is{metonymy} \textsc{covered for equipped} and \textsc{covered for adorned}, which both imply ends, respectively, something for which the patient/theme is equipped and someone for whose experience the patient/theme is adorned. If we assume that \iwi{PIE}{*\textit{kr̥h\textsubscript{2}-u-dʰh\textsubscript{1}-ó}-} ‘provided with cover, covered’ and its reflex \iwi{PIE}{*\textit{kruh\textsubscript{2}(-)dʰh\textsubscript{1}-ó}-} could develop senses of generic end-orientation and even end-orientation in relation to an experiencer considered capable of assigning positive value to the stimulus (= patient/theme), it would be conceivable that \mbox{\iwi{PIE}{*\textit{kruh\textsubscript{2}(-)dʰh\textsubscript{1}-ó}-}} expressed, on the one hand, the senses ‘equipped, prepared, ready’, and, on the other, ‘adorned, embellished, beautiful’. In this connection, it is appropriate to observe that Latin \iwi{Latin}{\textit{parātus}} ‘ready’ (: \iwi{Latin}{\textit{parāre}} ‘prepare’) and \iwi{Latin}{\textit{ōrnātus}} ‘adorned’ (: \iwi{Latin}{\textit{ōrnāre}} ‘fit out’), the \isi{resultative} adjectives corresponding to these two sense-domains, have been lexicalized as adjectives bestriding the Activity\is{activity} and Property Conceptual\is{property concept (PC)} domains. Correspondingly, we may suppose that \iwi{PIE}{*\textit{kruh\textsubscript{2}(-)dʰh\textsubscript{1}-ó}-} meant something like ‘equipped, adorned’.

\subsection{*\textit{kruh\textsubscript{2}-dʰh\textsubscript{1}-ó}- → *\textit{kréu̯h\textsubscript{2}(-)dʰh\textsubscript{1}-o/es}-}

The interaction of this quasi-Property Conceptual\is{property concept (PC)} meaning and the behavior of the proto-language with respect to word-formation is now of great importance. Like its sister, the merely \isi{resultative} \iwi{PIE}{*\textit{kruh\textsubscript{2}bʰ-ó}-} ‘heaped, covered’, \iwi{PIE}{*\textit{kruh\textsubscript{2}(-)dʰh\textsubscript{1}-ó}-} may be assumed to have been at some point analyzed as a simple thematic adjective \iwi{PIE}{*\textit{kruh\textsubscript{2}dʰh\textsubscript{1}-ó}-}. Given, however, its developed and thus distinctive PC\is{property concept (PC)} meaning, it is conceivable that speakers felt a greater need for that meaning to be expressed by means of an associated substantive. For the substantival expression of both Activity\is{activity}  Concepts and Property Concepts,\is{property concept (PC)} the readily available nominal stem formant was of course \iwi{PIE}{*-\textit{o/es}-}, the addition of which was prototypically accompanied by the installation of radical \textit{é}-grade for both primary\is{derivation!deradical}  and secondary derivation\is{derivation!stem-based} (e.g., \iwi{PIE}{*\textit{men}-} ‘think, feel’ $\rightarrow$ \iwi{PIE}{*\textit{mén-o/es}-} ‘mind, feeling’ > Ved.\ \iwi{Sanskrit}{\textit{mánas}-} ‘spirit’; \iwi{PIE}{*\textit{sru-tó}-} ‘flowing’ $\rightarrow$ \iwi{PIE}{*\textit{sréu̯-t-o/es}-} [i.e., *\textit{sréu̯-to/e-s}-] ‘flow’ > Ved.\ \iwi{Sanskrit}{\textit{srótas}-} n.\ ‘stream’). It is therefore reasonable to suppose that our \iwi{PIE}{*\textit{kruh\textsubscript{2}dʰh\textsubscript{1}-ó}-} ‘decked, equipped, adorned’ was at some point the base of an \textit{s}-stem substantive \mbox{\iwi{PIE}{*\textit{kréu̯h\textsubscript{2}(-)dʰh\textsubscript{1}-o/es}-}} that expressed the \isi{resultative} state of being equipped or adorned as a \textit{nomen qualitatis} (cf.\ PDE \textit{prepared-ness} ‘readiness’).

\subsection{*\textit{χrusti}- < *\textit{krudʰ-s-ti}-}


Before considering the influential role of this \textit{s}-stem in our account, we should of course adduce evidence that such a stem existed in the first place. That evidence, while indirect, is all the same fairly clear. Old High German \iwi{OHG}{\textit{hrust}} ‘armor’ and OE \iwi{OE}{\textit{hyrst}} ‘decoration, armor’, being exact cognates, point to a \iwi{PGmc}{*\textit{χrusti}-} derived from \iwi{PGmc}{*\textit{χreuð}-} by means of a suffix which \citet[275]{Seebold1970} represents as *-\textit{sti-} (*\textit{hrud-sti-} $\rightarrow$ \iwi{OHG}{*\textit{hrust-i-z}} [f.] ‘Rüstung’). At face value, such a sequence is necessary to account for the -\textit{st}- cluster actually attested, as we would expect an underlying Pre-Proto-Germanic \iwi{PIE}{*\textit{krudʰ(h\textsubscript{1})-ti}-} to be reflected as \iwi{PGmc}{*\textit{χrussi}-} (cf.\ OE \iwi{OE}{\textit{\.gewiss}} ‘sure, certain’ : \iwi{PIE}{*\textit{u̯id-tó}-} ‘seen, known’). Moreover, though it is imaginable that \iwi{PGmc}{*\textit{χrussi}-} was refashioned as \iwi{PGmc}{*\textit{χrusti}-} or that the reflex of an original *\textit{krudʰ-ti}-\iw{PIE}{*\textit{krudʰ(h\textsubscript{1})-ti}-} developed an excrescent -\textit{s}- resulting in \iwi{PGmc}{*\textit{χrusti}-} as well, such claims would be ad hoc. It is therefore preferable to suppose that the preform was a \iwi{PIE}{*\textit{krudʰ-s-ti}-} of some description.

\subsection{*\textit{χrausta}- < *\textit{krou̯dʰ-s-to}- : *\textit{krudʰ-s-tó}-}

Since \iwi{PIE}{*-\textit{ti}-} in this non-primary context points to derivation from an \isi{adnominal} in *-\textit{tó}- \iw{PIE}{*-\textit{to}-} \citep[390]{Schindler1980}, a safe assumption would be that \iwi{PIE}{*\textit{krudʰ-s-ti}-} was derived from a \iwi{PIE}{*\textit{krudʰ-s-tó}-} (cf.\ Pre-Hitt. \iwi{Hitt}{*\textit{dalugaš}-} ‘length’ $\rightarrow$ \iwi{Hitt}{*\textit{dalugašta}-} ‘lengthy, long’ $\rightarrow$ \iwi{Hitt}{\textit{dalugašti}-} ‘lengthiness, length’, \citealt[148]{Jasanoff2004}). Indirect evidence for this \iwi{PIE}{*\textit{krudʰ-s-tó}-} may be derived from Old Norse \iwi{ON}{\textit{hraustr}} ‘strong, valiant, doughty’, which certainly reflects a \iwi{PGmc}{*\textit{χrausta}-} related to \iwi{PGmc}{*\textit{χreuðaną}} \citep[174]{Hill2003}, the lexical semantic basis being the connection between strength or vitality and preparedness, as evident in the relation between \iwi{German}{\textit{rüstig}} ‘healthy’ and \iwi{German}{\textit{rüsten}} ‘equip’, the reflex of a denominal\is{denominal!verb}  to \iwi{PGmc}{*\textit{χrusti}-} (cf.\ OE \iwi{OE}{\textit{hyrstan}} ‘adorn’). Formally, \iwi{PGmc}{*\textit{χrausta}-} ‘ready’ would thus reflect a \iwi{PIE}{*\textit{krou̯dʰ-s-to}-} interpretable as an adjective formed by insertion of \textit{o}-grade into the root of \iwi{PIE}{*\textit{krudʰ-s-tó}-} (cf.\ \iwi{PIE}{*\textit{bʰoi̯d-ro}-} ‘cutting’ [Goth. \iwi{Goth}{\textit{baitrs}} ‘bitter’] : \iwi{PIE}{*\textit{bʰid-ró}-} ‘id.’ [OE \iwi{OE}{\textit{bitter}}, etc.]).

\subsection{*\textit{kréu̯dʰh\textsubscript{1}-o/es}- ‘equipment, adornment’}

If by comparison we consider, for example, the relation between (1) the thematic adjective basic to the preform of OE \iwi{OE}{\textit{āst}} f.\ ‘kiln’ (< PGmc.\ \iwi{PGmc}{*\textit{aistō}-}) and (2) the \textit{s}-stem that is ancestral to Vedic \iwi{Sanskrit}{\textit{édhas}-} ‘fuel’ and basic to that thematic adjective (\iwi{PIE}{*\textit{h\textsubscript{2}ei̯dʰ-s-to}-} ‘hot’ $\leftarrow$ \iwi{PIE}{*\textit{h\textsubscript{2}ei̯dʰ-es}-} ‘heat’, \citealt[527]{Nussbaum1998}), it is appropriate to assert that the absolute zero-grade \iwi{PIE}{*\textit{krudʰ-s-tó}-} ‘equipped, adorned’ (< \iwi{PIE}{*\textit{krudʰh\textsubscript{1}-s-tó}-})\footnote{For loss of *\textit{h\textsubscript{1}} between obstruents, see \citealt[77]{Jasanoff2003}.} was derived from an \textit{s}-stem \iwi{PIE}{*\textit{kréu̯dʰh\textsubscript{1}-o/es}-} ‘state of being equipped/adorned; equipment, adornment’, whose preform *\textit{kréu̯h\textsubscript{2}dʰh\textsubscript{1}-o/es}-\iw{PIE}{*\textit{kréu̯h\textsubscript{2}(-)dʰh\textsubscript{1}-o/es}-} reconstructed above had presumably undergone \citegen{Hackstein2002} rule of laryngeal deletion (*CH.CC > *C.CC). The question that now confronts us is the relation between this \iwi{PIE}{*\textit{kréu̯dʰh\textsubscript{1}-o/es}-} and the preform of Proto-Germanic \iwi{PGmc}{*\textit{χreuðaną}}.

\subsection{πλήθω (re-)revisited}

This question is best addressed by considering the history of \iwi{AG}{πλήθω} ‘be/become full’, which \citetalias[ 482--483]{Rix2001} classifies as the reflex of a present in \iwi{PIE}{*-\textit{dʰe/o}-},\is{*-\textit{dʰe/o}-present} namely \iwi{PIE}{*\textit{pléh\textsubscript{1}-dʰe/o}-}, also reflected in the diathetically divergent \iwi{Avestan}{\textit{frāda}-} ‘increase (tr.\ act./ intr.\ mid.)’ of Young and Old Avestan.\ What needs to be fundamentally explained in an account of the origin of this suffixal material is not this divergence in relation to some original semantic value of that material, but how the preform of πληθ- \iw{AG}{πλήθω}  and \textit{frād-} \iw{Avestan}{\textit{frāda}-} came to behave derivationally as a root when we know it is not a root at the deepest reconstructible stratum. The primary root was of course \iwi{PIE}{*\textit{pleh\textsubscript{1}}-} ‘fill’, which, like \iwi{PIE}{*\textit{\kinvbreve el}-} ‘cover’ or \iwi{PIE}{*\textit{kreh\textsubscript{2}}-} ‘cover, heap’, corresponded to an Activity\is{activity} Concept, rather than a Property Concept\is{property concept (PC)} (cf.\ \citealt[208--209]{Nussbaum2022}). As such, \iwi{PIE}{*\textit{pleh\textsubscript{1}}-} was eligible to form a root noun\is{root!noun} \iwi{PIE}{*\textit{pléh\textsubscript{1}-}/\textit{pl̥h\textsubscript{1}}-´} with the value of a \textit{nomen actionis} ‘filling’. Since, moreover, the specific lexical meaning of \iwi{PIE}{*\textit{pleh\textsubscript{1}}-} was of such a kind that the predication of \iwi{PIE}{*\textit{pl̥h\textsubscript{1}-éh\textsubscript{1}}} ‘with filling’ could express a state of being filled and thus full (cf.\ ‘with covering’ ${\approx}$ ‘covered’), the \isi{stative} expression could be used in the \isi{causative} \isi{periphrasis} \iwi{PIE}{*\textit{pl̥h\textsubscript{1}-éh\textsubscript{1}}} \textit{dʰeh\textsubscript{1}}-\iw{PIE}{*\textit{dʰeh\textsubscript{1}}-} ‘provide with filling, get \textit{x} full’, on the basis of which was formed a compound\is{compounding} \iw{PIE}{*\textit{pl̥h\textsubscript{1}(-)dʰh\textsubscript{1}-ó}-} *\textit{pl̥h\textsubscript{1}-dʰh\textsubscript{1}-ó}- ‘filling; filled, full’. Like the \isi{resultative} adjective \iwi{PIE}{*\textit{pl̥h\textsubscript{1}-nó}-} ‘filled, full’ (Ved.\ \iwi{Sanskrit}{\textit{pūrṇá}-} ‘full’) beside \iwi{PIE}{*\textit{pélh\textsubscript{1}-n-o/es}-} ‘fullness’ (: \iwi{Sanskrit}{\textit{parīṇas}-} ‘fullness’, Av.\ \textit{parənah-vant-} \iw{Avestan}{\textit{parənahvant}-} ‘copious’, \citealt[132, 174]{Rau2009}), *\textit{pl̥h\textsubscript{1}-dʰh\textsubscript{1}-ó}- \iw{PIE}{*\textit{pl̥h\textsubscript{1}(-)dʰh\textsubscript{1}-ó}-} ‘filled, full’, indirect evidence for which may be derived from the reflexes of nominals belonging to clearly de-thematic categories (e.g., \textit{ā}-stem Locrian \iwi{AG}{πληθα} ‘assembly, the commons’ < \iwi{PIE}{*\textit{pléh\textsubscript{1}-dʰh\textsubscript{1}-e-h\textsubscript{2}}-} ‘the full, fullness’, \textit{u}-stem Lat. \iwi{Latin}{\textit{plēbēs}} ‘the commons’ {<}{<} \iwi{PIE}{*\textit{pléh\textsubscript{1}-dʰh\textsubscript{1}-u}-}; v.\  \citealt{Nussbaum2014b}), appears to have been the base of an \textit{s}-stem *\textit{pléh\textsubscript{1}-dʰh\textsubscript{1}-o/es}- \iw{PIE}{*\textit{pléh\textsubscript{1}(-)dʰh\textsubscript{1}-o/es}-}  ‘fullness’ ancestral to \iwi{AG}{πλῆθoς} (cf.\ \iwi{Avestan}{*\textit{frādah}-} in Av.\ personal name \textit{Daŋ́hu-frādah-} \iw{Avestan}{\textit{Daŋ́hufrādah}-} ‘whose increase is of the land’, Yt.13.116, \citealt[682]{Bartholomae1904}).

\subsection{*\textit{pléh\textsubscript{1}-dʰh\textsubscript{1}-o/es}- ‘fullness’ : *\textit{pléh\textsubscript{1}-dʰh\textsubscript{1}-e/o}- ‘get full’}

Since that \textit{s}-stem was the substantival exponent of a Property Concept\is{property concept (PC)} \textsc{full}, it could serve as a means of derivational inspiration on the model of (primary) adjectival roots\is{root!adjectival} interpretable as basic to \textit{s}-stems and to prototypically radical \textit{é}-grade simple thematic presents \is{thematic!present} (e.g., \iwi{PIE}{*\textit{tep}-} ‘hot’ $\rightarrow$ \iwi{PIE}{*\textit{tép-o/es}-} ‘heat’ [> Ved.\ \iwi{Sanskrit}{\textit{tápas}-} ‘id.’] : \iwi{PIE}{*\textit{tép-e/o}-} ‘be/become hot’ [Ved.\ \textit{tápa-ti/te} \iw{Sanskrit}{\textit{tápati/te}} ‘is/becomes hot’], \iwi{PIE}{*\textit{g\textsuperscript{u̯h}er}-} ‘warm’ $\rightarrow$ \iwi{PIE}{\textit{*g\textsuperscript{u̯h}ér-o/es}-} ‘warmth’ [\iwi{AG}{θέρoς} ‘summer’] : \iwi{PIE}{\textit{*g\textsuperscript{u̯h}ér-e/o}-} ‘(be) warm’ [θέρoμαι\iw{AG}{θέρω, θέρομαι} ‘become warm’], etc.). While the ultimate origin of the latter category is mysterious,\footnote{On the \isi{thematization} of the \isi{proto-middle}, see \citealt{Jasanoff2023}.} Rau (\citeyear[263--265]{Rau2013}; this volume) has observed what appears to be an archaic correlation of active morphology and \isi{intransitive} semantics, which may point to the frequent -- though not necessarily exclusive -- \isi{anticausative} use of the \isi{proto-middle}. There may have been a degree of lability\is{labile verbs} comparable to the (anti)\isi{causative} \is{anticausative} use of PDE \textit{get} (e.g., \textit{get} (x) \textit{warm}; cf.\ PDE \textit{become} : NHG \iwi{German}{\textit{bekommen}} ‘get’). In any case, given the cooccurrence of such Calandish derivatives,\is{Caland!system (CS)}  namely radical \textit{é}-grade \textit{s}-stem substantives and largely \isi{anticausative} radical \textit{é}-grade simple thematic presents,\is{thematic!present} it is highly likely that  \iw{PIE}{*\textit{pléh\textsubscript{1}(-)dʰh\textsubscript{1}-o/es}-} *\textit{pléh\textsubscript{1}-dʰh\textsubscript{1}-o/es}- ‘fullness’ invited the suffixally substitutive formation of a *\textit{pléh\textsubscript{1}-dʰh\textsubscript{1}-e/o}-\iw{PIE}{*\textit{pléh\textsubscript{1}(-)dʰh\textsubscript{1}-e/o}-} ‘get full’, whose Greek reflex is active and \isi{intransitive} and whose Iranian one either happened to be active and \isi{transitive} or was an original \isi{anticausative} active remodeled as a middle corresponding to a \isi{causative} active backformed\is{backformation} from that middle in a manner suggested by \citet[155]{Rau2009}. However it might have been, the fundamental idea is that, for the purpose of derivation, the sequence \iw{PIE}{*\textit{pleh\textsubscript{1}(-)dʰh\textsubscript{1}}-}  *\textit{pleh\textsubscript{1}dʰh\textsubscript{1}}- was treated as a root with the meaning ‘filled, full’.

\subsection{*-\textit{dʰh\textsubscript{1}e/o}- and *-\textit{éh\textsubscript{1}i̯e/o}- compared}

Such an account is not designed to deny the existence of \isi{*-\textit{dʰe/o}-present}s \iw{PIE}{*-\textit{dʰe/o}-} with \isi{intransitive} semantics. Indeed, Greek evidence makes it very clear that such a formant existed (\iwi{AG}{φλεγέθω} ‘burn’, \iwi{AG}{φθινύθω} ‘perish’, \iwi{AG}{βαρύθω} ‘am heavy’, \iwi{AG}{φαέθω} ‘shine’, \iwi{AG}{μινύθω} ‘diminish’; see \citealt[268]{Magni2008}; \citealt[18--24]{RothsteinDowden2022}). What may be suggested concerning the origin of the category of \isi{*-\textit{dʰe/o}-present}s is an inference based on cases like \iw{PIE}{*\textit{pléh\textsubscript{1}(-)dʰh\textsubscript{1}-e/o}-}  *\textit{pléh\textsubscript{1}-dʰh\textsubscript{1}-e/o}- ‘get full’, where the “root-based” formation\is{derivation!deradical} from a base ending in \iwi{PIE}{*-\textit{dʰh\textsubscript{1}}-} was still analyzable with a morphological boundary before \iwi{PIE}{*-\textit{dʰh\textsubscript{1}}-}. Thus, a reanalyzed \iwi{PIE}{*-\textit{dʰh\textsubscript{1}e/o}-} could even be used to form primary verbal\is{derivation!deradical} stems with \isi{intransitive} meaning (e.g., \iwi{PIE}{*\textit{i̯eh\textsubscript{2}}-} : OCS \iwi{OCS}{\textit{jadǫ}} ‘go’ : Lith.\ \textit{jóju}, \textit{jóti}\iw{Lith}{\textit{jóti}, 1sg. prs. \textit{jóju}} ‘id.’, \iwi{PIE}{*\textit{pelh\textsubscript{2}}-} : \iwi{AG}{πελάθω/πλᾱ́θω} ‘approach’ : \iwi{AG}{πελάζω} ‘id.’). 
In this respect, a largely \isi{anticausative} \iwi{PIE}{*-\textit{dʰh\textsubscript{1}e/o}-} would, in terms of the history of the \isi{possessive} \isi{adnominal} system, intimately resemble the \isi{stative} *-\textit{éh\textsubscript{1}i̯e/o}-,\iw{PIE}{*-\textit{eh\textsubscript{1}-i̯e/o}-} which originated in the addition of \iwi{PIE}{*-\textit{i̯e/o}-} to the instrumental of root noun\is{root!noun} abstracts used in nominal sentences (\citealt[147]{Jasanoff2004}). Since these instrumental root nouns\is{root!noun} were usually predicated of themes or non-agents generally (e.g., \mbox{*\textit{pédom}} \textit{tr̥séh\textsubscript{1}}  \iw{PIE}{*\textit{tr̥séh\textsubscript{1}}} ‘the ground is dry’, *\textit{h\textsubscript{3}rḗg̑ōn h\textsubscript{2}n̥réh\textsubscript{1}} \iw{PIE}{*\textit{h\textsubscript{2}n̥réh\textsubscript{1}}} ‘the king is strong’, etc.),\footnote{I use “theme” to refer an argument of which a property is predicated.} the historically complex suffix has a \isi{stative} or \isi{anticausative} value. Though in principle a *\textit{h\textsubscript{1}rudʰ-éh\textsubscript{1}-i̯e/o}-\iw{PIE}{*\textit{h\textsubscript{1}rudʰ-eh\textsubscript{1}-i̯e/o}-} might have been a \isi{causative} (i.e., \isi{factitive}) ‘redden (\textit{x})’ like other adnominally \is{adnominal} based \isi{*-\textit{i̯e/o}-present}s, such as what are arguably the preforms of \iwi{AG}{κρύπτω} and \iwi{AG}{πλάσσω}, the fact that it meant ‘be/become red’ is plausibly attributed to the use of \iwi{PIE}{*\textit{h\textsubscript{1}rudʰ-éh\textsubscript{1}}} ‘with redness’ in \isi{stative} predication of themes.\footnote{For a formal analysis of the development of *-\textit{eh\textsubscript{1}}-\iw{PIE}{*-\textit{éh\textsubscript{1}-/-h\textsubscript{1}}-} into a verbal stem formant, see \citealt[15--23]{Grestenberger2023}.} Like hypostatic *-\textit{ó}-\iw{PIE}{*-\textit{o}-} and \iwi{PIE}{*-\textit{ii̯o}-} in nominals, \iwi{PIE}{*-\textit{i̯e/o}-} could, as one of its functions, merely render an already inflected substantive -- and thus uninflectable \isi{adnominal} -- inflectable. Once a unitary *-\textit{éh\textsubscript{1}i̯e/o}-\iw{PIE}{*-\textit{eh\textsubscript{1}-i̯e/o}-} ‘be/become (\textit{x})’ emerged, it could be used in primary derivation\is{derivation!deradical} (e.g., Lat. \iwi{Latin}{\textit{iaceō}} ‘lie’ : \iwi{Latin}{\textit{iaciō}} ‘throw’). A comparable development would be the case for \iwi{PIE}{*-\textit{dʰh\textsubscript{1}e/o}-}, whose denominal \is{denominal!verb} origin from compounds\is{compounding} based on the \isi{causative} \iwi{PIE}{*\textit{dʰeh\textsubscript{1}}-} \isi{periphrasis} would follow from the greater importance and thus frequency of mentioning agents in causation, and whose largely \isi{intransitive} semantics would be a consequence of the tendency for verbal adjectives\is{verbal adjective} to group patiency and resultativity\is{resultative} together for the sake of characterization \citep{Haspelmath1994}. Because verbal adjectives\is{verbal adjective} in \iwi{PIE}{*-\textit{dʰh\textsubscript{1}-ó}-} were largely statives\is{stative} of patient/theme-oriented and \isi{resultative} sense (*\textit{pl̥h\textsubscript{1}-dʰh\textsubscript{1}-ó}- ‘filled’),\iw{PIE}{*\textit{pl̥h\textsubscript{1}(-)dʰh\textsubscript{1}-ó}-} the adjectival abstracts formed from them on the basis of their \isi{resultative} \isi{stative} meaning (or secondary simple \isi{stative} sense) could inspire the creation of largely \isi{intransitive}, radical \textit{é}-grade “simple” thematic presents \is{thematic!present} formed from roots (or derivationally root-like bases) basic to Caland systems.\is{Caland!system (CS)}  Accordingly, the largely \isi{intransitive} value of the suffix \iwi{PIE}{*-\textit{dʰh\textsubscript{1}e/o}-} would be an inheritance from the formation of which its immediate ancestor *\textit{R(é)(-)dʰh\textsubscript{1}-e/o-} was a subtype: radical \textit{é}-grade simple thematic presents \is{thematic!present} (*\textit{R(é)-e/o-}).

\subsection{*\textit{kréu̯dʰh\textsubscript{1}-e/o}- ‘render equipped, adorned’}


For our purpose, the question is how the \iwi{PIE}{*\textit{pléh\textsubscript{1}(-)dʰh\textsubscript{1}-e/o}-} ancestral to  \iwi{PIE}{*\textit{pléh\textsubscript{1}-dʰh\textsubscript{1}e/o}-} was formed. By appeal to an \textit{s}-stem \iwi{PIE}{*\textit{pléh\textsubscript{1}(-)dʰh\textsubscript{1}-o/es}-} ‘fullness’ interpretable as derivative of a root-like base \iwi{PIE}{*\textit{pleh\textsubscript{1}(-)dʰh\textsubscript{1}}-} of Property Conceptual\is{property concept (PC)} sense, the \textit{s}-stem could easily have inspired a verbal formation in \iwi{PIE}{*\textit{-e/o}-} (e.g., \iwi{PIE}{*\textit{g\textsuperscript{u̯h}er}-} ‘warm’ $\rightarrow$ \iwi{PIE}{\textit{*g\textsuperscript{u̯h}ér-o/es}-} ‘warmth’ : \iwi{PIE}{\textit{*g\textsuperscript{u̯h}ér-e/o}-} ‘(be) warm’ :: \iwi{PIE}{*\textit{pléh\textsubscript{1}(-)dʰh\textsubscript{1}-o/es}-} ‘fullness’ : \textit{x} = \iwi{PIE}{*\textit{pléh\textsubscript{1}(-)dʰh\textsubscript{1}-e/o}-}; thus \iwi{PIE}{*\textit{pleh\textsubscript{1}(-)dʰh\textsubscript{1}}-} ‘full’ $\rightarrow$ \iwi{PIE}{*\textit{pléh\textsubscript{1}(-)dʰh\textsubscript{1}-e/o}-} ‘get full’). In the same way, I argue, our \iwi{PIE}{*\textit{kréu̯dʰh\textsubscript{1}-o/es}-} ‘state of being equipped, adorned, equipment, adornment’, interpretable as derivative of a quasi-Property Conceptual\is{property concept (PC)} \iwi{PIE}{*\textit{kreu̯dʰh\textsubscript{1}}-} ‘equipped, adorned’, inspired the formation of a \iwi{PIE}{*\textit{kréu̯dʰh\textsubscript{1}-e/o}-} ‘render equipped, adorned’ whose transitivity\is{transitive} would be the result of whatever underlies the transitivity\is{transitive} of the Avestan reflex of \iwi{PIE}{*\textit{pléh\textsubscript{1}(-)dʰh\textsubscript{1}-e/o}-} ‘get full’ (> \iwi{Avestan}{\textit{frāda}-} ‘increase [tr.\ act./intr.\ mid.]’) and whose Proto-Germanic reflex \iwi{PGmc}{*\textit{χreuðaną}} (> OE \iwi{OE}{\textit{hréodan}*} ‘deck, adorn’) consequently came to mean ‘equip, adorn’.

\section{Conclusion}

The main purpose of this contribution has been the argument that, most notably among a group of formally and semantically similar verbal derivatives,\is{derivation!verbal} Greek \iwi{AG}{κρύπτω} ‘cover, hide, bury’ and Old English \iwi{OE}{\textit{hréodan}*} ‘deck, adorn’ are not only related in terms of root-etymology, but also paradigmatic of the two ways in which a verb can be de-adjectival: stem-derived\is{derivation!stem-based}  (\iwi{AG}{κρύπτω}) vs.\ root-derived\is{derivation!deradical} \mbox{(\iwi{OE}{\textit{hréodan}*}).} While I have in general focused on the morphological and derivational semantic foundations underlying the formation of the etyma of both verbs and their relatives, it has been my particular intention to describe the conceptual and lexical semantic basis for a type of root-based derivation\is{derivation!deradical} that in IE terms may be designated as “de-adjectival”, namely Caland system verbal\is{derivation!verbal} derivation\is{Caland!system (CS)}  (see the introduction to this volume). The formal prerequisite to this description has been the identification of the origin of the roots ancestral to those of \iwi{AG}{κρύπτω} and \iwi{OE}{\textit{hréodan}*} in nominal derivation and the predicative\is{predicative!instrumental} instrumental (anti)\isi{causative} \is{anticausative} \isi{periphrasis} -- a type of construction closely, but not exclusively, associated with the Caland system\is{Caland!system (CS)}  and the actional expression of Property Concepts.\is{property concept (PC)} 

Despite the subsequent importance of Property Conceptuality\is{property concept (PC)} in this contribution, I have argued that the primary root underlying all the forms compared was in lexical semantic terms a verbal one \iwi{PIE}{*\textit{kreh\textsubscript{2}}-} ‘heap, strew, cover’ (Latvian \textit{krãju}, \textit{krât}\iw{Latv}{\textit{krât, krãju}} ‘put aside, accumulate, amass’) corresponding to an Activity\is{activity} Concept and basic to a \textit{u}-stem verbal noun \iwi{PIE}{*\textit{kró/éh\textsubscript{2}-u}-} ‘heap, covering’ whose instrumental was employed with \isi{anticausative} \iwi{PIE}{*\textit{bʰuh\textsubscript{x}}-} ‘become’ and \isi{causative} \iwi{PIE}{*\textit{dʰeh\textsubscript{1}}-} ‘put, make’. While the expression of an action involving this substantive as a possessum was in the present system handled by means of a derivative in \iwi{PIE}{*-\textit{i̯e/o}-} \is{*-\textit{i̯e/o}-present} (\iwi{PIE}{*\textit{kreh\textsubscript{2}-u-i̯e/o}-} ‘provide with covering’), whose \isi{reanalysis} as a primary verb\is{derivation!deradical} gave rise to a historically secondary verbal\is{root!secondary}  \iwi{PIE}{*\textit{kreu̯h\textsubscript{2}}-} (OCS \textit{kryti}\iw{OCS}{\textit{kryti, kryjǫ}} ‘cover’), constructions with the light verbs \iwi{PIE}{*\textit{dʰeh\textsubscript{1}}-} and \iwi{PIE}{*\textit{bʰuh\textsubscript{x}}-} were required elsewhere. I have argued that \iwi{PIE}{*\textit{kr̥h\textsubscript{2}-u-bʰ(u̯)-ó}-} ‘heaped, covered’ ({>}{>} \iwi{PIE}{*\textit{kruh\textsubscript{2}(-)bʰ-ó}-}), a thematic compound\is{compounding} formed on the basis of the latter construction, gave rise to the secondary\is{root!secondary} verbal root \iwi{PIE}{*\textit{kreu̯bʰ}-} (Tocharian B \iwi{TB}{\textit{kraup}-} ‘gather, amass’) and that the origin of this root lies in stem-based de-adjectival derivation.\is{derivation!stem-based} In a manner comparable to the use of \iwi{PIE}{*-\textit{i̯e/o}-} \is{*-\textit{i̯e/o}-present} to represent the functionally \isi{adnominal} \iwi{PIE}{*\textit{kréh\textsubscript{2}-u-h\textsubscript{1}}} ‘with covering’ in the derivative \iwi{PIE}{*\textit{kreh\textsubscript{2}-u-i̯e/o}-} ‘provide with covering’, the inflectable \isi{adnominal} \iwi{PIE}{*\textit{kruh\textsubscript{2}(-)bʰ-ó}-} ‘covered’ was the base of a \iwi{PIE}{*\textit{kruh\textsubscript{2}bʰ-i̯e/o}-} ‘render covered, heap (upon)’ whose reflex \iwi{PIE}{*\textit{krubʰ-i̯e/o}-} was the preform \iwi{AG}{κρύπτω}. 

The crucial contrast in terms of de-adjectival derivation is with the preform of \iwi{OE}{\textit{hréodan}*} ‘deck, adorn’, which arguably originated in a similar way, namely, from a compound\is{compounding} \iwi{PIE}{*\textit{kr̥h\textsubscript{2}-u-dʰh\textsubscript{1}-ó}-} ‘covered’ ({>}{>} \iwi{PIE}{*\textit{kruh\textsubscript{2}(-)dʰh\textsubscript{1}-ó}-}) formed on the basis of the \isi{causative} \isi{periphrasis} \iwi{PIE}{*\textit{kréh\textsubscript{2}-u-h\textsubscript{1}}} \textit{dʰeh\textsubscript{1}}-\iw{PIE}{*\textit{dʰeh\textsubscript{1}}-} ‘provide with covering’, but whose ancestral root’s meaning emerged as a result of a lexical semantic development that endowed the \isi{resultative} \isi{verbal adjective} \iwi{PIE}{*\textit{kruh\textsubscript{2}(-)dʰh\textsubscript{1}-ó}-} ‘covered’ with the capacity to express an idea in the vicinity of a Property Concept,\is{property concept (PC)} i.e., ‘equipped, ready’ or ‘adorned’. This adjacency to Property Conceptuality\is{property concept (PC)} was the precondition for the emergence of Calandish derivational behavior.\is{Caland!system (CS)}  Once \mbox{\iwi{PIE}{*\textit{kruh\textsubscript{2}(-)dʰh\textsubscript{1}-ó}-}} ‘equipped, adorned’ became the base of an \textit{s}-stem substantive \iwi{PIE}{*\textit{kréu̯h\textsubscript{2}(-)dʰh\textsubscript{1}-o/es}-}  ‘state of being equipped/adorned’, the resulting \iwi{PIE}{*\textit{kréu̯dʰh\textsubscript{1}-o/es}-} ‘equipment, adornment’, now analyzed as a primary derivative,\is{derivation!deradical} implied the existence of a derivative belonging to a verbal stem type closely associated with the Caland system:\is{Caland!system (CS)}  the radical \textit{é}-grade simple thematic present\is{thematic!present} (\iwi{PIE}{*\textit{kréu̯dʰh\textsubscript{1}-e/o}-} ‘render equipped/adorned; provide with equipment/adornment’). In sum, the preforms of both Greek \iwi{AG}{κρύπτω} ‘cover’ and Proto-Germanic \iwi{PGmc}{*\textit{χreuðaną}} ‘deck, adorn’ not only exemplify the complexity of de-adjectival derivation, but also offer some insight into the ways in which morphology and derivational semantics interact with our conceptual world.

\section*{Abbreviations}
\begin{tabularx}{.48\textwidth}{lQ}
act.\ & active\\
Alb.\ & Albanian\\
Av.\ & Avestan\\
f.\ & feminine\\
Hitt. & Hittite\\
Hom. & Homeric\\
Hsch. & Hesychius\\
IE & Indo-European\\
intr.\ & intransitive\\
Lat.\ & Latin\\
Latv.\ & Latvian\\
Lith.\  & Lithuanian\\
m.\ & masculine\\ 
mid.\ & middle\\
n.\ & neuter\\
NHG & New High German\\
OCS & Old Church Slavonic\\
\end{tabularx}
\begin{tabularx}{.48\textwidth}{lQ}
OE & Old English\\
OCS & Old Church Slavonic\\
OE & Old English\\
OHG & Old High German\\
OIr.\ & Old Irish\\
OLith.\  & Old Lithuanian\\
ON & Old Norse\\
PC & Property Concept\\
PDE & Present-day English\\ 
PGmc.\ & Proto-Germanic\\
PGk.\ & Proto-Greek\\
PIE & Proto-Indo-European\\
TA & Tocharian A\\
TB & Tocharian B\\
tr. & transitive\\
Ved.\ & Vedic\\
Yt.\ & Yašt\\
\end{tabularx}

\sloppy\printbibliography[heading=subbibliography,notkeyword=this]
\end{document} 
